\documentclass[11pt]{article}

\usepackage[margin=1in]{geometry}
\usepackage{amsmath}
\usepackage{amssymb}
\usepackage{mathtools}
\usepackage{enumitem}

% ----- light notation helpers -----
\newcommand{\Lang}{\theta_{L}}          % left leg angle
\newcommand{\Rang}{\theta_{R}}          % right leg angle
\newcommand{\Lvel}{\omega_{L}}          % left leg angular velocity
\newcommand{\Rvel}{\omega_{R}}          % right leg angular velocity
\newcommand{\ileg}{d}                   % interleg signal
\newcommand{\degs}{^{\circ}}            % degree symbol (\degs is reserved in LaTeX)
\newcommand{\sgn}{\operatorname{sgn}}
\newcommand{\ptp}{\operatorname{ptp}}
\newcommand{\corr}{\rho}
\DeclareMathOperator*{\argmax}{arg\,max}

\title{S3 Physics: The Model-Free View of a Gait Window}
\author{h-care pipeline, stage S3}
\date{}

\begin{document}
\maketitle

\begin{abstract}
Stage S3 turns a window of two-leg inertial data into a small set of physical
claims about the body, computed with no fitted parameters. The centerpiece is the
\emph{swap rule}: a walk/stand call built from the geometry of the legs alternating,
not from a trained boundary. Around it sit four \emph{anchors} (periodicity,
antiphase, gravitational stability, gyroscopic energy), each of which is then put
through a rate-invariance audit that keeps only the anchors describing the body and
discards the ones that secretly describe the sampling grid. This note states the
physics and the exact quantities the code computes.
\end{abstract}

\section{Signals and windowing}

Each recording provides four low-pass-filtered channels sampled at a canonical rate
$f_s = 100~\text{Hz}$:
\[
\Lang(t),\quad \Rang(t) \quad[\degs], \qquad
\Lvel(t),\quad \Rvel(t) \quad[\degs/\text{s}],
\]
the left/right leg tilt angles and their angular velocities. Analysis proceeds on
non-overlapping windows of length $T = 2~\text{s}$, i.e. $N = \lfloor T f_s \rceil = 200$
samples, and windows never straddle a recording gap.

The single most informative construction is the \textbf{interleg signal}, the
difference of the two leg angles:
\begin{equation}
\ileg(t) \;=\; \Lang(t) - \Rang(t).
\label{eq:interleg}
\end{equation}
When the legs swing in opposition (walking), $\ileg$ oscillates through zero with
large amplitude. When the body is still (standing), $\ileg$ sits near a constant.
The plausible cadence band for this population is
$[\,f_\text{lo}, f_\text{hi}\,] = [\,0.13, 3.0\,]~\text{Hz}$. Since the interleg signal
completes one cycle per stride and a stride is two steps, that is about 16 to 360 steps
per minute, a deliberately permissive bound. Every lag-based quantity below is restricted
to periods inside this band.

\section{The swap rule: a zero-parameter walk/stand call}

Walking is the legs \emph{alternating}. The swap rule measures exactly that: how many
times the interleg signal commits to one leg leading, then commits to the other. It
uses one physically-motivated constant, the commit threshold
\begin{equation}
\delta = 1.0\degs,
\end{equation}
set at or just above five times the standing noise: five times the raw standing noise
measured $0.76$ to $0.88\degs$ across files (so the raw floor is about $0.15$ to
$0.18\degs$), and $\delta$ rounds that up to $1.0\degs$. It is a noise gate, not a tuned
decision boundary: a commit must clear real motion, not jitter.

\subsection{Committed crossings with hysteresis}

Scan $\ileg$ sample by sample and maintain a committed side
$s \in \{+1, -1, 0\}$ (start uncommitted, $s=0$). A new commit is recorded only when
the signal clears the floor \emph{on the opposite side} to the last commit:
\begin{equation}
\text{record } +1 \text{ if } \ileg > +\delta \text{ and } s \neq +1;
\qquad
\text{record } -1 \text{ if } \ileg < -\delta \text{ and } s \neq -1,
\end{equation}
updating $s$ to the recorded side each time. Small wiggles that do not cross
$\pm\delta$, and repeat excursions to the side already held, change nothing: this is
hysteresis. By construction the recorded commits strictly alternate in sign, so the
number of alternations (the \textbf{swap count}) is
\begin{equation}
K \;=\; \max\!\big(0,\; m - 1\big),
\end{equation}
where $m$ is the number of recorded commits.

\subsection{The discrete verdict}

The swap count maps to a three-way verdict:
\begin{equation}
\text{verdict}(K) =
\begin{cases}
\textsc{standing}, & K = 0,\\[2pt]
\textsc{ambiguous}, & K = 1,\\[2pt]
\textsc{walking}, & K \geq 2.
\end{cases}
\label{eq:verdict}
\end{equation}
The $K=1$ case is genuinely ambiguous rather than a fudge factor: one leg passing the
other happens both when you take a step and when you shift your weight. Two full
alternations ($K \geq 2$) is the smallest unambiguous signature of a stride.

\section{The per-file rest zero}

The raw interleg signal carries a constant per-subject offset from sensor mounting and
zeroing. If that offset is large, $\ileg$ can stay above $-\delta$ through real gait,
so the swap rule never commits to the second side and misreads walking as standing.
The fix is to recenter $\ileg$ on the value it holds at rest.

Let $c$ be the \textbf{rest zero}, estimated as the median of $\ileg$ over a still span,
and used only inside the swap count:
\begin{equation}
K \;=\; \text{swap\_count}\big(\ileg - c\big).
\end{equation}
The posture descriptors of Section~\ref{sec:anchors} deliberately stay on raw $\ileg$,
because there the offset \emph{is} the posture.

\paragraph{Choosing the span, label-free.} A span qualifies as rest by the swap rule's
own standing verdict, applied on its locally centered copy:
\begin{equation}
\text{is\_rest}(\text{span}) \iff
\text{swap\_count}\big(\text{span} - \operatorname{median}(\text{span})\big) = 0.
\end{equation}
The estimator then chooses:
\begin{enumerate}[leftmargin=2em, itemsep=2pt]
  \item the opening $3~\text{s}$ if it is rest (recordings usually begin at rest); else
  \item the median of the stillest true-rest span anywhere in the file, ranked by the
        smallest peak-to-peak interleg range $\ptp(\ileg)$ (offset-free, so it ranks
        stillness, not posture); else
  \item the whole-file median, returned flagged as \emph{untrusted} (the recording
        never rests).
\end{enumerate}
Reusing the validated standing verdict rather than inventing a new stillness threshold
keeps this step parameter-free.

\section{Stride-adaptive analysis span}

A fixed $2~\text{s}$ window holds fewer than two strides once the stride period
approaches $2~\text{s}$, so slow gait reads $K \in \{0,1\}$ and is under-called. Since
the deployment population (assistive and rehab gait) lives in exactly this slow regime,
the swap count is evaluated over a span sized to the local stride period.

\subsection{Stride period by autocorrelation}

For the centered interleg segment $x = \ileg - \operatorname{mean}$ of length $n$ (the
same signal the swap rule reads; an in-phase non-gait motion leaves $\ileg$ flat and so
yields no spurious period to grow the window on), form the normalized autocorrelation at
lag $k$:
\begin{equation}
r(k) \;=\; \frac{\displaystyle\sum_{i} x_i\, x_{i+k}}{\displaystyle\sum_{i} x_i^2},
\qquad r(0) = 1.
\end{equation}
Restrict the lag to the cadence band,
$k \in [\,k_\text{lo}, k_\text{hi}\,]$ with
$k_\text{lo} = \lfloor f_s / f_\text{hi} \rceil$ and
$k_\text{hi} = \lfloor f_s / f_\text{lo} \rceil$. The stride period is the first
\emph{local maximum after the autocorrelation first dips below zero}:
\begin{equation}
k^{\star} = \argmax_{k \,\geq\, k_0} r(k),
\qquad
k_0 = \min\{\,k : r(k) < 0\,\},
\end{equation}
accepted only if $r(k^{\star}) \geq 0.35$ (a rhythm strong enough to be gait). Skipping
the lag-0 shoulder is essential: a plain $\argmax$ over all lags grabs the
monotonic-decay shoulder at the shortest lag and reports the floor period, not the true
one. If the ACF never dips (no full cycle resolved) or the peak is too weak, no period
is returned and the base window is kept (standing has no detectable period).

\subsection{Sizing the swap window}

Let the base half-window be $N/2$ and the cap be
$H_\text{max} = \lfloor H_{\max,s}\, f_s / 2 \rceil$ with $H_{\max,s} = 6~\text{s}$. Given a
detected period $k^{\star}$, the half-span used for the swap count at a cell center is
\begin{equation}
h =
\begin{cases}
N/2, & \text{no period detected (standing),}\\[2pt]
\operatorname{clip}\big(k^{\star},\; N/2,\; H_\text{max}\big), & \text{periodic (about two strides).}
\end{cases}
\end{equation}
The verdict is then \eqref{eq:verdict} evaluated on $\ileg$ over
$[\,c_\text{cell}-h,\; c_\text{cell}+h\,]$. Standing windows are never lengthened, so a
still bout is not stretched into its walking neighbors.

\section{The grow fallback}

Some slow, large-amplitude strides defeat the period detector: it returns no period,
the window stays at $2~\text{s}$, holds a single hump, and reads $K=1$. As a rescue,
whenever the adaptive verdict is \emph{not} walking, the swap count is recomputed over
the full $6~\text{s}$ cap and promoted to \textsc{walking} only if that span
\emph{genuinely alternates}. Over halves $A$ (first) and $B$ (second) of the max span,
all three conditions must hold:
\begin{equation}
\min\big(\ptp(d_A),\, \ptp(d_B)\big) \geq 8\degs,
\qquad
-\corr(\Lang, \Rang) \geq 0.5,
\qquad
K_{\text{max span}} \geq 2.
\end{equation}
The peak-to-peak gate requires a real swing on both sides of the window (not jitter or a
one-sided excursion). The anti-correlation gate is the key discriminator: two stray
standing weight-shifts can also produce $K=2$ over $6~\text{s}$, but their legs are not
antiphase, so requiring leg anti-correlation prevents merging them into a false walk.
This is a geometric test of alternation, not a second period detector, so it sidesteps
the failure mode it is rescuing.

\section{The four anchors}
\label{sec:anchors}

An anchor is a claim about the body. S3 reports four, each a scalar per window, and each
required to carry a rate-invariance verdict (Section~\ref{sec:audit}).

\paragraph{Antiphase.} The legs swing in opposition during gait. With Pearson
correlation $\corr$,
\begin{equation}
\text{antiphase} \;=\; -\,\corr\big(\Lang, \Rang\big) \in [-1, 1],
\end{equation}
signed so that larger means more antiphase. This is \emph{necessary but not sufficient}
for walking: gait requires the legs to alternate, so walking is antiphase; but non-gait
alternating movements (weight shifts, rocking) are antiphase too, so antiphase alone
cannot confirm gait. (In-phase motions such as a ramp or sit-to-stand are positively
correlated, so a low or negative value here rules gait out.)

\paragraph{Gravitational stability.} Steadiness of the gravity-referenced tilt. With
$\sigma_L, \sigma_R$ the standard deviations of the two leg angles in the window,
\begin{equation}
\text{grav\_stab} \;=\; \frac{1}{1 + \tfrac{1}{2}\big(\sigma_L + \sigma_R\big)} \in (0, 1].
\end{equation}
It approaches $1$ when posture holds still (standing) and falls toward $0$ when the tilt
sweeps (walking).

\paragraph{Periodicity.} Rhythm \emph{strength}, amplitude-independent. Using the same
normalized autocorrelation $r(k)$ per leg, take the tallest peak at a nonzero lag whose
period is a plausible stride and short enough that at least two cycles fit in the window
($k \leq n/2$):
\begin{equation}
p(x) \;=\; \operatorname{clip}\!\Big(\max_{k_\text{lo} \leq k \leq n/2} r(k),\; 0,\; 1\Big),
\qquad
\text{periodicity} \;=\; \tfrac{1}{2}\big(p(\Lang) + p(\Rang)\big).
\end{equation}
Here $1.0$ is perfectly periodic and $\approx 0$ is no rhythm. It is a within-window
property (steady rhythm \emph{here}), and is known to drop at cadence changes rather than
at true arrhythmia.

The $k \leq n/2$ bound has a consequence worth stating: at the fixed $2~\text{s}$ window
only rhythm at $\geq 1~\text{Hz}$ (stride period $\leq 1~\text{s}$) fits twice, and the slow
assistive/rehab strides (down to $0.33~\text{Hz}$) that are exactly the deployment population
do not. With fewer than two cycles in the window the $\max$ over $[k_\text{lo}, n/2]$ lands on
the short-lag autocorrelation \emph{shoulder} rather than a resolved stride peak, so the anchor
turns into a smoothness proxy that cannot tell rhythmic slow gait from a non-rhythmic smooth
sweep: a $0.4~\text{Hz}$ stride and a plain linear ramp both read $\approx 0.5$ in a
$2~\text{s}$ window. It is not near-zero, it is \emph{uninformative} there. So alongside this
anchor S3 also
reports \textbf{adaptive periodicity}, the same quantity measured over the stride-adaptive
span the swap count already sizes to about two strides (up to the $6~\text{s}$ cap), which
resolves rhythm down to $\approx 0.33~\text{Hz}$.
For standing and non-periodic cells that span is the base window, so the two agree; it
widens only where a stride is too slow to fit twice in $2~\text{s}$. Crucially, the
\emph{audited} anchor stays the fixed-window one: the rate-invariance verdict of
Section~\ref{sec:audit} and every downstream consumer read the fixed-window
\text{periodicity}, and the adaptive version rides alongside as a descriptor, exactly as
the stride-adaptive swap verdict rides beside the fixed one.

\paragraph{Gyroscopic energy.} Total rotational energy in the window,
\emph{summed} over samples:
\begin{equation}
\text{gyro\_energy} \;=\; \sum_{i} \big(\Lvel\big)_i^{\,2} \;+\; \sum_{i} \big(\Rvel\big)_i^{\,2}
\qquad [\text{deg}^2/\text{s}^2].
\end{equation}
Because it is a sum rather than a mean, it scales with the sample count $N$, so it is
really a claim about the sampling grid, not the body. It is defined this way on purpose:
the audit below is expected to reject it, re-deriving a verdict this quantity has failed
before.

\section{The rate-invariance audit}
\label{sec:audit}

An anchor that describes the body should not change when the same bout is measured at a
different sample rate. The audit tests this directly: decimate each window by a factor of
two through the anti-aliasing FIR filter (a genuine bandwidth cut from
$100~\text{Hz}$ to $50~\text{Hz}$, zero-phase so peaks do not move), recompute each
anchor, and measure the change.

\paragraph{Change on the right scale.} Each anchor is compared on its natural scale.
Bounded scores and correlations use an absolute delta; ratio-scale magnitudes use a
floored relative delta:
\begin{equation}
\Delta_a =
\begin{cases}
\big|\,a_\text{native} - a_\text{dec}\,\big|, & a \in \{\text{periodicity}, \text{antiphase}, \text{grav\_stab}\}\ (\text{absolute}),\\[6pt]
\dfrac{\big|\,a_\text{native} - a_\text{dec}\,\big|}{\big|\,a_\text{native}\,\big| + \varepsilon_a}, & a = \text{gyro\_energy}\ (\text{relative},\ \varepsilon_a = 1.0).
\end{cases}
\end{equation}
Mixing the two would distort the result: framing a bounded score of $0.60$ vs $0.55$ as a
relative change exaggerates it near zero, while framing a genuine energy shift as an
absolute one ignores that its size depends on the bout amplitude.

\paragraph{Gated to the motion regime.} The audit is computed only on windows the swap
rule calls \textsc{walking}. These four anchors are descriptors of locomotion; a standing
window has no cadence and no rhythm, so comparing that silence to its decimated self
measures decimation noise, not rate dependence, and would falsely condemn a good anchor.
Gating on the independently validated swap rule keeps the audit label-free.

\paragraph{Verdict.} Collect $\Delta_a$ over all walking windows and compare the median to
a tolerance $\tau = 0.10$:
\begin{equation}
\text{verdict}(a) =
\begin{cases}
\text{invariant}, & \operatorname{median}(\Delta_a) \leq \tau,\\[2pt]
\text{rate\_dependent}, & \operatorname{median}(\Delta_a) > \tau.
\end{cases}
\end{equation}
The bounded anchors pass; gyroscopic energy fails, as intended, because its value tracks
the number of samples rather than the moving body. A rate-dependent verdict is a finding,
not an error: it tells any downstream reader (human or agent) that the anchor may be
discussed but must never carry a load-bearing claim about the body.

\section{The agent layer: from figures to hypotheses}

Everything above is deterministic: the code computes the anchor table, runs the
audit, and draws the figures with no model in the loop. A separate language-model
agent then sits on top of that output. Its job is the one thing code cannot do, look
at the physics and say what it means, but under strict provenance rules that keep it
honest.

\paragraph{The figures are the only evidence.} The agent cannot see the corpus. It is
given read-only access to the rendered figures and the summary tables, and nothing
else. Each per-trial figure stacks three panels on one time axis in seconds:
\begin{enumerate}[leftmargin=2em, itemsep=2pt]
  \item the two leg angles $\Lang, \Rang$, with the \emph{human label} drawn as a
        background band;
  \item the interleg swap signal $\ileg = \Lang - \Rang$ with the $\pm\delta$ commit
        bands and the stride-adaptive swap verdict as colored markers along the top;
  \item the anchor timeline (antiphase, periodicity, gravitational stability).
\end{enumerate}
The human label is drawn as the \emph{background}, not as ground truth. Where the
physics track disagrees with the label band is exactly what the agent is asked to
explain: slow walking still labeled as walking, a split-stance bout labeled standing,
a moving ramp labeled standing.

\paragraph{Reading order is precomputed.} To avoid reading every figure blind, the
deterministic core ranks each trial by physics-vs-label disagreement. For a trial it
reports the fraction of walk-labeled windows the swap rule does not call walking and
the fraction of stand-labeled windows it does, and it histograms the physical
\emph{cause} of each conflict (legs in phase, a single un-alternated hump, a
sub-threshold swing, or a standing stretch the physics reads as gait). The agent is
handed this ranking and the anchors' rate-invariance verdicts, and reads the audit
figure and the top-conflict trials first.

\paragraph{Output: falsifiable, window-anchored hypotheses.} The agent does not label
windows and does not train. It performs rule \emph{discovery}, emitting a set of
hypotheses about the body and the labels. Each hypothesis is a structured record with
a one-sentence statement, a confidence level, the anchors it relies on, a caveat
stating what would falsify it, and, crucially, a list of \emph{evidence windows}, each
naming a recording, a trial, a real time interval, and the anchors it points at.

\paragraph{A provenance gate re-checks every hypothesis in code.} After the agent
replies, deterministic validation enforces two hard rules, and any hypothesis that
breaks either is flagged rather than silently kept:
\begin{enumerate}[leftmargin=2em, itemsep=2pt]
  \item \textbf{Window-level provenance.} Every hypothesis must point at real windows
        with a valid time interval ($t_\text{end} > t_\text{start}$) and at least one
        named anchor. A general claim with no window is rejected.
  \item \textbf{Respect the audit.} Every anchor the hypothesis leans on is
        cross-referenced against its rate-invariance verdict. A claim resting on a
        \emph{rate\_dependent} anchor (such as gyroscopic energy) is permitted only if
        it explicitly acknowledges that verdict; otherwise it is flagged.
\end{enumerate}
Validated hypotheses are written one per line, each annotated with the rate-invariance
verdict of its anchors and an ok-or-flagged status; the raw agent reply is kept
alongside for audit. A hypothesis whose physics contradicts a region's human label
sets a label-audit flag, routing that window to human review.

In one line: the deterministic core turns physics into figures, the agent turns
figures into falsifiable window-anchored hypotheses, and code re-checks their
provenance before any hypothesis is recorded.

\end{document}
